% Full finite decision, information-comparison and transport proofs.

\begin{annexitem}{annex:p17.information-policy-value}{thm:p17-information-policy-value}{Signalwise optimization and zero-probability outputs}
\begin{proof}
For each signal define the finite real numbers
\[
                         b_z(a)=\sum_s\pi_s p_{sz}u(a,s).
\]
Regrouping the finite sum gives
\(J=\sum_z\sum_a\delta(a\mid z)b_z(a)\).
For each \(z\), the inner sum is a convex combination of the \(b_z(a)\),
so it is at most their maximum. Since \(A\) is finite and nonempty, one
of its actions attains that maximum. Choosing one maximizing action at
each of the finitely many signals gives a deterministic policy attaining
the sum of the maxima. This proves existence and the first value formula.

If \(q_z>0\), division by \(q_z\) identifies \(b_z(a)\) with
\(q_z\sum_s\pi_s^z u(a,s)\), and multiplication by this positive number
commutes with taking the maximum. If \(q_z=0\), every nonnegative term
\(\pi_s p_{sz}\) in its sum is zero. Thus \(b_z(a)=0\) for every action.
The entire contribution of that signal is zero for every policy, without
assigning a posterior on a null event. These observations prove the second
formula and all policy assertions.
\end{proof}
\end{annexitem}

\begin{annexitem}{annex:p17.information-value-bounds}{thm:p17-information-value-bounds}{Ignoring a signal and comparison with statewise choice}
\begin{proof}
Let \(a_0\) attain the finite no-information maximum. The policy that
chooses \(a_0\) at every signal has expected payoff
\[
 \sum_{s,z}\pi_s p_{sz}u(a_0,s)
 =\sum_s\pi_s u(a_0,s)=V_0,
\]
using each row sum of \(P\). It is feasible, so \(V\geq V_0\).
For any policy and state,
\[
 \sum_{z,a}p_{sz}\delta(a\mid z)u(a,s)
 \leq \max_a u(a,s)\sum_{z,a}p_{sz}\delta(a\mid z)
 =\max_a u(a,s).
\]
Multiply by \(\pi_s\), sum over states and maximize over policies.
This gives \(V\leq V_{\mathrm{PI}}\). The identity experiment attains
the upper bound by selecting a maximizing action in each observed state.

Subtracting \(V_0\) gives the two information-value bounds.
A fixed payoff-unit cost subtracts the same \(c\) from every policy's
payoff, so the new optimum is \(V-c\). Its comparison with \(V_0\)
is equivalent to \(c\leq V-V_0\). Feasibility of ignoring the signal
was used in the first inequality; no conclusion for an altered action
set or an unavoidable state-dependent cost is inferred.
\end{proof}
\end{annexitem}

\begin{annexitem}{annex:p17.garbling-simulation}{thm:p17-garbling-simulation}{Statewise policy simulation and the experiment preorder}
\begin{proof}
Each entry of \(\widehat\delta\) is nonnegative, and
\[
 \sum_a\widehat\delta(a\mid z)
   =\sum_y M_{zy}\sum_a\delta(a\mid y)=1.
\]
Thus it is a policy. Conditional on any fixed state \(s\), its action
probabilities are
\[
 \sum_z p_{sz}\widehat\delta(a\mid z)
   =\sum_{z,y}p_{sz}M_{zy}\delta(a\mid y)
   =\sum_y q_{sy}\delta(a\mid y).
\]
This equality is independent of the prior and proves identical statewise
payoffs for every action/payoff array. It also proves the value inequality
by simulating an optimal \(Q\)-policy.

The identity matrix is row stochastic and gives reflexivity. Products
of row-stochastic matrices are row stochastic: their entries are
nonnegative and summing a product row first over its last index gives
the first factor's row sum. If \(PM=Q\) and \(QN=R\), then
\(P(MN)=R\), proving transitivity.
Mutual garbling is reflexive and symmetric by definition, and transitive
by composing both directional transformations.

The identity state experiment transforms into \(P\) using the matrix
\(P\) itself. Sending every signal to one output gives the single-signal
experiment. To see that equivalence need not preserve the number of
signals, compare the single-signal experiment with an experiment whose
rows all equal one fixed probability vector on two signals. A constant
randomized output generates the latter from the former, and forgetting
the output gives the reverse transformation.
\end{proof}
\end{annexitem}

\begin{annexitem}{annex:p17.finite-separation-minimax}{thm:p17-finite-separation-minimax}{Compactness, projection and a finite minimax proof}
\begin{proof}
\paragraph{Finite-dimensional compactness used here.}
A bounded sequence of real numbers has a convergent subsequence:
successively bisect a closed interval containing the sequence and retain
a half containing infinitely many of its terms. The nested intervals
have lengths tending to zero and a common point by completeness of the
reals. Choose successively increasing indices from the retained infinite
sets; the resulting subsequence converges to that point.
For a bounded sequence of vectors, perform this selection successively
for its finitely many coordinates. The final subsequence converges in
every coordinate. If the vectors lie in a closed set, their limit lies
in that set.

Consequently a continuous real function on a nonempty closed bounded
set is bounded and attains its extrema. Unboundedness would supply a
sequence of values increasing beyond every integer, while a convergent
subsequence of their arguments would force a finite limiting value.
A sequence of values approaching the finite infimum likewise has a
convergent argument subsequence, and continuity attains the infimum.
The supremum argument is the same. These are ordinary classical
finite-dimensional real arguments; no effective algorithm on unspecified
real inputs is asserted.

Finite products of closed bounded sets are closed and bounded, and the
same subsequence argument shows that a continuous image of such a set
is closed and bounded. For closedness, take a convergent sequence in
the image, choose preimages, and use a convergent subsequence in the
domain. For boundedness, an unbounded image sequence would give the
same contradiction to continuity. In particular the simplex and a
finite convex hull, which is a linear image of a simplex, have these
properties.

\paragraph{Separation by a closest point.}
The continuous function \(d^2(b,c)\) attains a minimum at \(c_0\in C\).
For any \(c\in C\) and \(0<t\leq1\), convexity puts
\(c_0+t(c-c_0)\) in \(C\). Minimality and direct expansion give
\[
 0\leq -2t\langle b-c_0,c-c_0\rangle+t^2d^2(c,c_0).
\]
Divide by \(t\) and let \(t\) decrease to zero. This yields
\(\langle h,c-c_0\rangle\leq0\).
Since \(b\notin C\), \(h\ne0\); hence
\[
 \langle h,b\rangle-\langle h,c_0\rangle
       =\langle h,h\rangle>0.
\]
This proves exactly the displayed separation, including its direction.

\paragraph{A finite matrix game.}
Write \(r_i\) for the row vectors of \(H\), and let \(C\) be their
convex hull in column-coordinate space. Put
\[
                        v=\max_{x\in C}\min_j x_j.
\]
The maximum exists by the compactness just proved.
Mixing the rows gives precisely \(C\), and minimizing a linear function
over the column simplex chooses a smallest coordinate. Thus \(v\)
is the left side of the minimax equality.

Fix \(\varepsilon>0\), and let
\(O_\varepsilon=\{y:y_j\geq v+\varepsilon\ \hbox{for every }j\}\).
No point of \(C\) lies in this orthant, by the definition of \(v\).
For \(x\in C\), its closest orthant point is the coordinatewise vector
\(y_j=\max(x_j,v+\varepsilon)\): each squared coordinate difference
is minimized by that choice. The resulting squared distance is a
continuous function of \(x\), so it has a minimum at \(x_\varepsilon\).
Let \(y_\varepsilon\) be its associated orthant point and
\(d=y_\varepsilon-x_\varepsilon\). Then \(d\) is nonnegative and
nonzero.

This pair minimizes squared distance over \(C\times O_\varepsilon\).
In particular, comparison with
\((x_\varepsilon+t(x-x_\varepsilon),y_\varepsilon)\) and the same
expansion as above give
\[
                    \langle d,x-x_\varepsilon\rangle\leq0
                    \qquad(x\in C).
\]
Whenever \(d_j>0\), the corresponding coordinate of \(y_\varepsilon\)
is \(v+\varepsilon\). Therefore
\[
 \langle d,x_\varepsilon\rangle
   =(v+\varepsilon)\sum_jd_j-\sum_jd_j^2
   <(v+\varepsilon)\sum_jd_j.
\]
Normalize by the strictly positive sum to obtain a column probability
vector \(\beta_\varepsilon=d/\sum_jd_j\).
Every row then satisfies \(r_i\cdot\beta_\varepsilon<v+\varepsilon\).
Take \(\varepsilon=1,1/2,1/3,\ldots\) and a convergent subsequence in
the closed bounded simplex. Its limit \(\beta_0\) satisfies
\(r_i\cdot\beta_0\leq v\) for every row.

Conversely, for any row mixture \(\alpha\) and column mixture \(\beta\),
\[
 \min_{\beta'}\alpha^{\mathsf T}H\beta'
       \leq\alpha^{\mathsf T}H\beta
       \leq\max_{\alpha'}(\alpha')^{\mathsf T}H\beta.
\]
Taking the relevant outer extrema gives the weak minimax inequality.
The constructed \(\beta_0\) gives the reverse inequality, so equality
holds. A maximizing row mixture \(\alpha_0\) exists because its row
payoff vector attains the definition of \(v\).
The two optimality inequalities show
\[
 \alpha^{\mathsf T}H\beta_0
 \leq \alpha_0^{\mathsf T}H\beta_0
 \leq \alpha_0^{\mathsf T}H\beta
\]
for every \(\alpha,\beta\), with middle value \(v\).
This proves attainment and the saddle-point assertion.
\end{proof}
\end{annexitem}

\begin{annexitem}{annex:p17.finite-blackwell}{thm:p17-finite-blackwell}{The finite comparison equivalences}
\begin{proof}
\paragraph{Garbling reproduces loss vectors.}
Suppose \(PM=Q\). Given choices \(a_y\in\mathcal A\) after \(Q\),
choose \(\widehat a_z=\sum_yM_{zy}a_y\) after \(P\).
Convexity puts \(\widehat a_z\) in \(\mathcal A\).
For every state, finite-sum regrouping gives
\[
 \sum_zp_{sz}(\widehat a_z)_s
       =\sum_yq_{sy}(a_y)_s.
\]
Thus \(B(Q,\mathcal A)\subseteq B(P,\mathcal A)\).
For clarity, the attainable sets themselves are nonempty, closed,
bounded and convex: the choices range over the finite product
\(\mathcal A^r\), and the risk-vector map is linear and continuous.
The preceding compactness argument gives closedness and boundedness;
convex combinations of choice arrays prove convexity.

\paragraph{Loss inclusion gives value comparison.}
Fix a finite action/payoff array \(u\). Associate each action with
the loss vector \((-u(a,s))_{s\in S}\), and let \(\mathcal A\)
be the convex hull of these finitely many vectors. Choosing a point
of this hull at a signal is exactly choosing a randomized action's
statewise loss vector; every such point has finite convex weights
on the original actions. Thus inclusion of attainable loss vectors
allows every \(Q\)-policy's statewise payoff vector to be attained
with \(P\). It preserves expected payoff for any prior, giving
condition 3. Condition 3 implies condition 4 by using its specified
full-support prior.

\paragraph{One full-support prior and all payoffs force a garbling.}
Let \(\mathcal M\) be the set of row-stochastic \(r\times t\) matrices,
and let
\[
                         C=\{PM:M\in\mathcal M\}.
\]
The set \(\mathcal M\) is a nonempty closed bounded convex set:
each entry lies in \([0,1]\), the row-sum equalities are closed,
and a constant row distribution supplies a member.
Its linear image \(C\) has the same closed, bounded and convex properties.

If \(Q\notin C\), the separation lemma gives an array \(H\) such that
\[
                 \langle H,Q\rangle>
                     \max_{M\in\mathcal M}\langle H,PM\rangle.
\]
Use the signal set \(Y\) as the action set and define the finite real
payoffs
\[
                            u(y,s)=H_{sy}/\pi_s^0.
\]
Every denominator is positive. A policy after \(P\) is exactly a
row-stochastic matrix \(M\) with columns indexed by these actions.
Its expected payoff is \(\langle H,PM\rangle\), so the maximum on
the right is \(V(\pi^0,P,u)\).
After \(Q\), choosing action \(y\) when signal \(y\) occurs gives
payoff \(\langle H,Q\rangle\). Therefore
\[
                       V(\pi^0,Q,u)>V(\pi^0,P,u),
\]
contradicting condition 4. Thus \(Q\in C\), which is condition 1.
All four conditions have now been linked in a cycle of implications.
The proof did not infer statewise comparisons from one fixed payoff
problem; it used the quantifier over every payoff array.
\end{proof}
\end{annexitem}

\begin{annexitem}{annex:p17.blackwell-trace-route}{thm:p17-blackwell-trace-route}{The trace condition and Blackwell's finite minimax route}
\begin{proof}
Given a real \(t\times n\) matrix \(D\), let \(\mathcal A\) be the
convex hull of its \(t\) rows. After observing \(y\) under \(Q\),
choose row \(y\) as the loss vector. Its risk vector is
\(\operatorname{diag}(QD)\). Since \(P\supset Q\), this vector
is achieved after \(P\) by some choices \(a_z\in\mathcal A\).
Write each \(a_z\) as a convex combination of the rows of \(D\),
with weights \(M_{zy}\). These weights form a row-stochastic \(M\).
Their risk vector is \(\operatorname{diag}(PMD)\), proving the
claimed equality of each diagonal entry and hence the traces.

For the weaker condition define
\[
                         h(D,M)=\operatorname{Trace}((Q-PM)D).
\]
This is bilinear in \(D,M\), apart from its term independent of \(M\),
which does not affect bilinearity after mixtures with weights summing
to one. More precisely it is affine in each variable separately and
preserves every convex combination in each variable.
The cube \([0,1]^{t\times n}\) is the convex hull of its finitely many
zero-one vertices. The explicit weight of a vertex is the product
of the selected coordinate values or their complements; these weights
sum to one and have the desired coordinate means.
The row-stochastic matrices are the convex hull of the finitely many
matrices selecting one output in each row. For a deterministic selection
\(f\), its weight is \(\prod_zM_{z,f(z)}\); finite distributivity
proves normalization and reproduces every entry of \(M\).

Apply the proved finite minimax theorem to the payoff matrix whose
rows are the cube vertices, columns are these deterministic matrices,
and entries are \(h(D,M)\). Mixtures reproduce the original continuous
cube and row-stochastic domains because \(h\) preserves those mixtures.
It follows that
\[
               \max_D\min_M h(D,M)=\min_M\max_D h(D,M),
\]
and the minimizing matrix \(M_0\) exists.
The hypothesis says \(\min_M h(D,M)\leq0\) for every \(D\), so
the displayed equality gives \(h(D,M_0)\leq0\) for every cube member.

Put \(U=Q-PM_0\). To inspect one entry \(U_{sy}\), choose \(D\)
with entry \(D_{ys}=1\) and every other entry zero.
Then \(\operatorname{Trace}(UD)=U_{sy}\leq0\).
Every row of \(U\) has sum zero, because both \(Q\) and \(PM_0\)
are row stochastic. A finite row of nonpositive numbers with sum
zero has every entry zero. Hence \(U=0\) and \(PM_0=Q\).
This proves the source route without assuming an unproved minimax
theorem or exchanging the two quantifiers directly.
\end{proof}
\end{annexitem}

\begin{annexitem}{annex:p17.posterior-splitting}{thm:p17-posterior-splitting}{Posterior means and an explicit implementing matrix}
\begin{proof}
For each state,
\[
 \sum_{z:q_z>0}q_z\pi_s^z
   =\sum_{z:q_z>0}\pi_s p_{sz}
   =\sum_z\pi_s p_{sz}=\pi_s.
\]
The middle equality holds because all terms in a zero column-mass
sum are nonnegative and hence individually zero.

Conversely, for \(\pi_s>0\), the proposed entries are nonnegative
and their row sum is
\(\sum_zq_z\nu_s^z/\pi_s=1\).
For \(\pi_s=0\), choose any row probability vector, for example
one concentrated on the first signal. If \(q_z>0\), the barycenter
equation at such a state forces \(\nu_s^z=0\), since its nonnegative
weighted sum is zero.
It follows for every \(s,z\) that
\[
                         \pi_s p_{sz}=q_z\nu_s^z.
\]
For \(\pi_s=0\) and \(q_z=0\), both sides are zero as well.
Summing over states gives exactly \(q_z\) as the signal mass.
At positive mass, division by \(q_z\) returns the specified posterior
\(\nu^z\). Zero-weight signals have no posterior claim and no
contribution to the law. This proves the construction for all priors.
\end{proof}
\end{annexitem}

\begin{annexitem}{annex:p17.finite-posterior-order}{thm:p17-finite-posterior-order}{Conditional posterior averaging and the converse}
\begin{proof}
Suppose \(Q=PM\). Couple the state and both signals with mass
\(\pi_s p_{sz}M_{zy}\). The joint mass of \(z,y\) is \(q_zM_{zy}\),
and its second marginal is \(\widetilde q_y\).
For \(\widetilde q_y>0\), the posterior of state \(s\) is
\[
 \frac{\pi_s q_{sy}}{\widetilde q_y}
 =\sum_{z:q_z>0}\frac{q_zM_{zy}}{\widetilde q_y}\,\pi_s^z.
\]
The weights sum to one. Finite convexity therefore gives
\[
 \phi(\widetilde\pi^y)
 \leq\sum_{z:q_z>0}\frac{q_zM_{zy}}{\widetilde q_y}\phi(\pi^z).
\]
The finite Jensen inequality used here follows directly by induction
from two-point convexity: remove zero weights, separate one remaining
weight, and normalize the rest unless that weight is one.
Multiply by \(\widetilde q_y\), sum over its positive-mass outputs
and regroup. Terms in a zero-mass output vanish, and the row sums
of \(M\) give the desired posterior-law inequality.

For the converse fix any finite action/payoff problem and set
\[
                        \phi_u(p)=\max_a\sum_s p_su(a,s).
\]
This function is convex: each linear value at a mixture is at most
the same mixture of the two maxima. It is continuous. Indeed, if
each coordinate of \(p-p'\) has absolute value at most \(\delta\),
then each linear value changes by at most
\(\delta\sum_s\operatorname{abs}(u(a,s))\); the maximum of these
finitely many bounds also bounds the change of the maximum.
The value formula gives
\[
              V(\pi,P,u)=\sum_zq_z\phi_u(\pi^z).
\]
Thus the assumed convex-function comparisons imply value dominance
at the fixed full-support prior for every finite payoff problem.
The finite Blackwell theorem gives \(P>Q\).
The same argument used only maxima of finitely many linear functions,
which proves the stronger sufficiency assertion.
\end{proof}
\end{annexitem}

\begin{annexitem}{annex:p17.finite-convex-transport}{thm:p17-finite-convex-transport}{An explicit simplex embedding and all null rows}
\begin{proof}
First assume a transformation \(T\) as stated. Finite Jensen gives,
for each \(y\in X\),
\[
                         \phi(y)\leq\sum_xT_{yx}\phi(x).
\]
Multiply by \(\nu_y\), sum and use its transformed marginal \(\mu\).
This proves the convex-order inequalities.

Conversely assume those inequalities. Both coordinate functions \(x_i\)
and their negatives are continuous and convex. Applying the two
inequalities gives equality of each coordinate mean. Denote the common
mean vector by \(c\).
Put \(n=d+1\) and choose the explicit constants
\[
 R=\max_{x\in X}\left(1+\sum_{i=1}^d
                       \operatorname{abs}(x_i-c_i)\right),\qquad
 \varepsilon=\frac{1}{2nR}.
\]
The finite nonempty set gives \(R\geq1\).
Define an affine map \(\iota\) on \(\mathbb R^d\) by
\[
 \begin{aligned}
 \iota_i(x)&=1/n+\varepsilon(x_i-c_i)\quad(1\leq i\leq d),\\
 \iota_n(x)&=1/n-\varepsilon\sum_{i=1}^d(x_i-c_i).
 \end{aligned}
\]
Its coordinates sum to one. On \(X\), every deviation from \(1/n\)
is at most \(1/(2n)\) in absolute value, so every coordinate is
strictly positive. Convexity gives the same membership in the
probability simplex on the convex hull of \(X\).
The map is injective, with inverse coordinate formula
\[
                      x_i=c_i+(\iota_i(x)-1/n)/\varepsilon.
\]
The \(\mu\)-mean and \(\nu\)-mean of \(\iota(x)\) are both the
uniform vector \((1/n,\ldots,1/n)\).

Construct experiments on \(n\) states, each with signal labels \(X\):
\[
                       P_{ix}=n\mu_x\iota_i(x),\qquad
                       Q_{ix}=n\nu_x\iota_i(x).
\]
Nonnegativity and the preceding means show that every row sums to one.
At the uniform prior, the respective signal masses are \(\mu_x,\nu_x\).
At each positive-mass signal the posterior vector is \(\iota(x)\).
For any continuous convex function \(\psi\) on the simplex,
\(\psi\circ\iota\) is continuous and convex on the convex hull of \(X\).
The assumed order therefore gives the posterior convex-order inequality
for \(P,Q\). The preceding theorem supplies a row-stochastic \(M\)
with \(PM=Q\).

Let \(\lambda_{xy}=\mu_xM_{xy}\).
Summing the equality \(PM=Q\) over state indices gives
\[
                   \sum_x\lambda_{xy}=\nu_y.
\]
The individual coordinate equations give
\[
                   \sum_x\lambda_{xy}\iota_i(x)
                              =\nu_y\iota_i(y).
\]
For \(\nu_y>0\), define \(T_{yx}=\lambda_{xy}/\nu_y\).
Its row sums to one, and the affine inverse of \(\iota\) gives
\(\sum_xT_{yx}x=y\). For \(\nu_y=0\), define the row as a point
mass at \(y\) itself. It has mean \(y\), as required even at a null row.
Moreover the nonnegative column identity above implies
\(\lambda_{xy}=0\) for every \(x\) when \(\nu_y=0\).
Consequently
\[
 \sum_y\nu_yT_{yx}
   =\sum_y\lambda_{xy}
   =\mu_x\sum_yM_{xy}
   =\mu_x.
\]
This proves both the target marginal and the mean-preserving property
with every zero-mass case explicitly accounted for.
\end{proof}
\end{annexitem}

\begin{annexitem}{annex:p17.information-example-checks}{thm:p17-information-example-checks}{The service calculation, incomparable partitions and a null state}
\begin{proof}
For the binary processing example, the first diagonal product entry is
\((4/5)(3/4)+(1/5)(1/4)=13/20\), and the first off-diagonal entry
is \((4/5)(1/4)+(1/5)(3/4)=7/20\). The second row gives the symmetric
values, so \(PM=Q\). Every row of \(M\) is a probability vector.
Under the uniform prior, either signal has mass \(1/2\). The posterior
probability of its matching state is \(4/5\) under \(P\) and \(13/20\)
under \(Q\); guessing that state is optimal. This gives the stated values.
A reverse garbling would imply the reverse value inequality in this same
problem, contradicting \(4/5>13/20\).

The expected operating payoff without a signal is \(-10/5=-2\),
whereas servicing gives \(-3\), so the stated no-information value follows.
The alert mass is
\[
            (1/5)(4/5)+(4/5)(1/10)=4/25+2/25=6/25.
\]
Its bad-state posterior is \((4/25)/(6/25)=2/3\).
The no-alert bad mass is \(1/25\) and its total mass is \(19/25\),
giving posterior \(1/19\).
After an alert, operating has payoff \(-20/3<-3\), so servicing is
optimal. Otherwise operating has payoff \(-10/19>-3\).
Thus
\[
                V=(6/25)(-3)+(19/25)(-10/19)=-28/25.
\]
The difference from \(-2\) is \(22/25\).
Perfect information services only the bad state, giving
\(V_{\mathrm{PI}}=-3/5\) and \(\operatorname{EVPI}=7/5\).
Subtracting \(1/2\) gives \(-81/50\), strictly larger than
\(-2=-100/50\).

For the three-state example, a correct binary answer pays one and an
incorrect answer pays zero. Experiment \(P\) answers the state-1 question
perfectly. Under \(Q\), signal 1 identifies state 2 and permits a correct
negative answer; its probability is \(1/3\). Signal 2 leaves states
1 and 3 equally likely, so every randomized answer has success at most
\(1/2\) there. Total optimal success is
\(1/3+(2/3)(1/2)=2/3\).
The symmetric argument reverses the values for the state-2 question.
There is also a direct matrix obstruction: rows 2 and 3 of \(P\)
are identical, so rows 2 and 3 of every \(PM\) are identical, whereas
the corresponding rows of \(Q\) differ. Conversely rows 1 and 3 of
\(Q\) are identical but those of \(P\) differ.

For the transport example, the zero row has mean
\((-1)/2+1/2=0\), and each endpoint row is its own point mass and has
that endpoint as mean. Weighting the rows by \(\nu\) selects only
the zero row, producing exactly the claimed \(\mu\). Two-point convexity
gives the displayed inequality. When the comparison measure is instead
concentrated at 1, its mean under the affine convex function \(x\) is one,
whereas the endpoint mixture has mean zero, so the required inequality
does fail.

Finally, at prior \((1,0)\), every signal generated from either experiment
has all its positive posterior mass on state 1. The optimum in every
payoff problem is simply \(\max_a u(a,1)\), so the values agree.
The single-signal experiment has two identical rows. Multiplication by
any row-stochastic postprocessor preserves their equality, and thus
cannot produce the two different rows of the identity experiment.
More generally, a row with \(\pi_s=0\) contributes zero to every joint
mass and every expected payoff at that prior. Fixed-prior comparisons
cannot distinguish changes confined to such rows.
\end{proof}
\end{annexitem}
